% !TeX root = ../初等数论笔记.tex
% 由主文档通过 \input 载入；请编译主文档。
\setcounter{chapter}{12} % 保留原书章号。
\chapter{素数分布初步（选学）}
\label{ch:prime-distribution}

本章先用整除与同余研究特殊数列，再讨论正项级数与渐近式；后两部分使用级数收敛和实函数极限的基本概念。

“素数有无穷多个”只说明没有最后一个素数，并不说明素数出现得有多频繁。要比较稀疏程度，可以数出某个区间的素数，也可以考察素数倒数的累积增长。本章把初等可证的结论和需要进一步理论的定理明确分开，使读者理解每种说法究竟提供多少信息。

\section{特殊等差数列中的素数}

\begin{proposition}
形如 $4k+3$ 的素数和形如 $4k+1$ 的素数都各有无穷多个。
\end{proposition}
\begin{proof}
若模 $4$ 余 $3$ 的素数只有 $p_1,\ldots,p_r$，令 $N=4p_1\cdots p_r-1$。$N$ 是模 $4$ 余 $3$ 的奇数，故其素因子中至少一个 $q$ 模 $4$ 余 $3$，否则素因子全部余 $1$，乘积也余 $1$。但每个 $p_i$ 都不能整除 $N$，矛盾。

若模 $4$ 余 $1$ 的素数只有 $q_1,\ldots,q_s$，令 $A=2q_1\cdots q_s$，取奇数 $A^2+1$ 的素因子 $q$。此时 $q\nmid A$ 且 $A^2\equiv-1\pmod q$，由补充公式知 $q\equiv1\pmod4$；它又不同于原有全部 $q_i$，矛盾。空列表按空乘积为 $1$ 处理，同样成立。
\end{proof}

更一般的\emph{狄利克雷素数定理（Dirichlet's theorem on primes in arithmetic progressions）}断言：若正整数 $a,m$ 互素，则 $a+km$（$k\ge0$）含无穷多个素数。互素条件不可删去，因为公共因子会整除数列中的每一项。这里仅陈述一般定理，其解析证明见\cite{kedlaya-dirichlet}。

例如 $6k+5$ 与 $6k+1$ 都包含无穷多个素数；但 $6k+3$ 中除 $3$ 本身外全是合数。若 $d=\gcd(a,m)>1$，数列中的素数项只能等于 $d$，且只有在 $d$ 本身为素数时才有这种可能，因此至多出现一个素数项。反过来，互素条件保证无穷多个素数，是一个远强于“没有固定公因子障碍”的定理，不能仅凭障碍消失就宣布结论成立。

\section{素数间隔与素数倒数和}

给定 $L\ge1$，令 $N=(L+1)!$，则 $N+2,N+3,\ldots,N+L+1$ 都是合数：第 $j$ 项被 $j$ 整除且严格大于 $j$。这构造了任意长的连续合数段，但没有断言空隙长度与其所处位置成正比。

\begin{theorem}[素数倒数和发散]
$\sum_{p\text{ 为素数}}1/p$ 发散。
\end{theorem}
\begin{proof}
若该正项级数收敛，可取 $P$ 使 $\sum_{p>P}1/p<1/2$。对任意 $X>P$，由有限乘积不等式 $\prod_i(1-u_i)\ge1-\sum_i u_i$（$0\le u_i\le1$，可归纳证明），得到
\[
 \prod_{P<p\le X}(1-1/p)>\frac12.
\]
另一方面，将有限多个收敛几何级数相乘，由唯一分解，
\[
 \sum_{n\le X}\frac1n
 \le\prod_{p\le X}(1-1/p)^{-1}
 <2\prod_{p\le P}(1-1/p)^{-1}.
\]
第一不等式成立，是因为每个 $n\le X$ 的素因子均不超过 $X$，所以其倒数必出现在乘积展开中。右端为固定常数，但调和级数发散：每个区间 $2^{j-1}<n\le2^j$ 的倒数和至少为 $1/2$。矛盾。
\end{proof}

伯特兰命题（Bertrand's postulate）进一步保证对每个整数 $n>1$，区间 $(n,2n)$ 内至少有一个素数。这里将它作为不证的选学结论，相关讨论与推广见\cite{sondow-bertrand}。它与任意长合数段并不矛盾：前者控制相对尺度，后者只要求绝对长度任意大。

\section{素数计数函数与数量级}

定义\emph{素数计数函数（prime-counting function）} $\pi(x)=\#\{p\le x:p\text{ 为素数}\}$（$x\ge2$）。记 $f(x)\sim g(x)$ 表示 $f(x)/g(x)\to1$。\emph{素数定理（prime number theorem）}给出
\begin{equation}
 \pi(x)\sim\frac{x}{\log x}\qquad(x\to\infty),
 \label{eq:prime-number-theorem}
\end{equation}
其中 $\log$ 为自然对数。该定理在此引用而不证明，进一步阅读见\cite{kedlaya-pnt}。

渐近等价描述相对误差趋于零，并不意味着两边之差趋于零，也不直接给出某个有限 $x$ 下的误差界。例如 $\pi(100)=25$，而 $100/\log100\approx21.7$。实际计数、可证明上下界与渐近式是不同层次的信息。

\begin{example}[如何正确使用渐近等价]
在允许引用素数定理的前提下，估计区间 $(x,2x]$ 中的素数个数。
\end{example}
\begin{solve}
不能把渐近符号当作精确等号直接相减，而应将每一项除以共同尺度 $x/\log x$。由素数定理，
\[
 \frac{\pi(2x)}{x/\log x}
 =\frac{\pi(2x)}{2x/\log(2x)}\frac{2\log x}{\log(2x)}\longrightarrow2,
 \qquad \frac{\pi(x)}{x/\log x}\longrightarrow1.
\]
两式相减得到 $\pi(2x)-\pi(x)\sim x/\log x$。这里两项的主导常数不同，相减后仍有非零主项；若主项完全抵消，仅凭原来的渐近等价通常不能确定差的数量级。
\end{solve}

从 $1,\ldots,N$ 等可能地选一个整数时，选到素数的概率恰为 $\pi(N)/N$，所以素数定理给出它渐近于 $1/\log N$。这里的随机性来自“等可能抽取”的实验，而不是整数本身随机地成为素数；该比例也不说明不同整数的素性事件彼此独立。

\section{习题}

本章区分无穷性、长空隙、倒数和与渐近计数四种信息。前几类结论不能自动互相推出；引用素数定理时，应明确使用的极限方向和共同尺度，而不将渐近式当作有限范围的精确公式。

\begin{enumerate}
 \item 构造连续五个合数，并验证每个数的一个真因子。
 \item 若只知道某一剩余类中有无穷多个素数，能否推出其中素数的倒数和发散？说明为何仅凭“无穷多个”不足以推断。
 \item 利用素数定理证明 $\pi(2x)-\pi(x)\sim x/\log x$；注明这里使用了一个未在正文证明的定理。
 \item 用筛法核对 $\pi(100)=25$，比较计数结果与式\eqref{eq:prime-number-theorem}的表达式。
 \item\label{ex:distribution-sparse} 给出一个无穷正整数集合，其倒数和收敛；用求和公式说明收敛，而不只写出集合名称。
 \item\label{ex:distribution-six-five} 仿照欧几里得式构造，证明模 $6$ 余 $5$ 的素数有无穷多个，不使用一般狄利克雷定理。
 \item\label{ex:distribution-density} 引用素数定理证明 $\pi(x)/x\to0$，并解释为什么这与素数倒数和发散不矛盾。
 \item\label{ex:distribution-asymptotic} 构造正函数 $f,g$ 使 $f(x)\sim g(x)$ 但 $f(x)-g(x)\to+\infty$，指出相对误差与绝对误差的区别。
\end{enumerate}

